% EP20--EP23. Written as an editable source-level proof, 19 October 2026.
% Dependencies: represented full exact calculus; (T),(P), realised lifted
% pair/triangle ages; joint exact quotient inheritance (A130/A131/A194).
% Audit status: same-assistant source-relative, NOT formalised/refereed.
\section{Exact trace geometry of safe unary towers}
\label{sec:ep-exact-traces}

Throughout this and the next two sections, \(M=M_0\) is one
\emph{original} transitive homogeneous complete finite named symmetric
binary source with AP and semantic homomorphic CCO.  An actual
\emph{safe tower}
\[
  \mathcal K_0\longleftarrow\mathcal K_1
  \longleftarrow\cdots\longleftarrow\mathcal K_m
  \tag{ET}\label{eq:ep-tower}
\]
means a finite succession of \emph{maximal-exact-complement,
full rigid-groupoid character} constructions, valid for
AP+homomorphic-CCO sources by
\cref{thm:ep-ap-marked-completion,cor:ep-ap-safe-splitting}.
The SAP-specific V27
\cref{thm:safe-exact-split} is a separate stronger source used
for the hidden Ramsey provider; it must not be imposed on
the AP-only exact quotients of the present induction.
Each \(\mathcal K_j\) is the age of an actual
homogeneous prefix \(M_j\); the reducts are identified with the
preceding homogeneous sources using reasonableness.  We distinguish
original old binary atom names from the unary types introduced in
the tower.  Every current same-sort relation is symmetric.

We use the following \emph{source contracts}, all in the represented,
equality-inclusive convention. Exact domains are finite nonempty
intersections of reachable endpoint profiles; they are closed under
converse, represented product and nonempty meets. On a safe character
step, write \(\mathsf A_{uv}\) for the mask of old atoms allowed
between children \(u,v\).
\begin{enumerate}[label=(\textup{\alph*}),leftmargin=*]
\item[(T)] Every pair of compatible allowed lifted atomic edges has
an allowed output atom of the required endpoint child types.
\item[(P)] If reachable old profiles \(P_i\) have nonempty
intersection and each meets \(\mathsf A_{uv}\), then
\((\bigcap_iP_i)\cap\mathsf A_{uv}\ne\varnothing\).
\item[(R)] A lifted complete pair/triangle is realised precisely
when its old atoms and all its child phase masks are compatible.
The full character construction is reasonable, reduct-surjective,
and has a homogeneous realised age.
\end{enumerate}
The historical SAP-source form of (T),(P) is reproduced
in the V27 appendix. Their AP-only derivation from the
represented exact calculus, including the relative-\(K\)
identity correction and \emph{actual} real marked AP/CCO
completion, is proved in the supplementary source module
\cref{lem:ep-ap-orbital-parity,lem:ep-ap-TP,thm:ep-ap-marked-completion}
earlier in Part~I. The quotient assertions below
additionally use the represented
joint exact-axis homogeneity and CCO theorem from the exact
geometry of Part~I.  We do not infer any of these contracts from
formal associativity of an unspecified colouring algebra.

For an old same- or mixed-typed relation \(D\), put
\[
   \operatorname{tr}_{uv}(D)
   =D\cap(u\times v)
\]
as an \emph{actual} relation on real points of the designated
current child sorts.  Its atomic palette is \(D\cap\mathsf A_{uv}\).
In a composition of traces the middle point must lie in the
displayed middle sort; this restriction is essential.

\begin{lemma}[Primality on nonempty exact intersections]
\label{lem:ep-exact-meet-prime}
If old exact domains \(D_1,\ldots,D_r\) have nonempty intersection
and all \(\operatorname{tr}_{uv}(D_i)\ne\varnothing\), then
\[
    \operatorname{tr}_{uv}\Bigl(\bigcap_iD_i\Bigr)
            \ne\varnothing .
\]
The assertion is \emph{not} made for an empty old intersection.
\end{lemma}
\begin{proof}
Write each \(D_i\) as a finite intersection of reachable profiles.
Each constituent profile meets the allowed mask because \(D_i\)
does. The complete flattened family has nonempty old intersection,
so apply (P). The typed top represents the empty defining family
and may be omitted.
\end{proof}

\begin{theorem}[Exact trace product]
\label{thm:ep-trace-product}
Let \(D:s\to t\) and \(E:t\to w\) be nonempty old exact domains.
For children \(u,v,z\) over \(s,t,w\) suppose both
\(\operatorname{tr}_{uv}(D)\) and
\(\operatorname{tr}_{vz}(E)\) are nonempty. Then
\[
 \boxed{\ 
 \operatorname{tr}_{uv}(D)\ ;
 \operatorname{tr}_{vz}(E)
   =\operatorname{tr}_{uz}(D;E).\
 }\tag{TP}\label{eq:ep-tp}
\]
This is equality of actual relations, not just of nonempty
endpoint colour sets.
\end{theorem}
\begin{proof}
One inclusion is immediate. For the other, take an actual
endpoint atom \(c\in D;E\) that is allowed between \(u,z\).
Peircean transfer shows
\(D\cap(\{c\};E^\smile)\ne\varnothing\).
The first factor \(D\) meets \(\mathsf A_{uv}\) by hypothesis.
So does the old exact domain
\(F=\{c\};E^\smile\): choose
\(b_0\in E\cap\mathsf A_{vz}\), and apply (T) to the allowed
fork \(c,b_0^\smile\). By
\cref{lem:ep-exact-meet-prime} choose
\[
 a\in D\cap(\{c\};E^\smile)\cap\mathsf A_{uv}.
\]
A second choice is needed; merely selecting \(a\) does not fix
the phase of the other edge.  Put
\(H=\{a\}^\smile;\{c\}\). By Peircean transfer,
\(E\cap H\ne\varnothing\). The domain \(E\) meets
\(\mathsf A_{vz}\) by hypothesis, and (T) on the fork
\(a^\smile,c\) makes \(H\) meet it.
A second application of
\cref{lem:ep-exact-meet-prime} gives
\[
 b\in E\cap(\{a\}^\smile;\{c\})\cap\mathsf A_{vz}.
\]
The old triangle \(a,b,c\) is realised and all three prescribed
lifted pair types are allowed. Contract (R) realises its lifted
triangle. Homogeneity aligns its \(c\)-edge with any prescribed
realised endpoint pair of child types \(u,z\), giving an actual
middle vertex of type \(v\). If an input or output is logical
identity, this is the corresponding equality/two-point
degeneracy, \emph{not} an artificial equality-coloured edge
joining distinct vertices.
\end{proof}

\begin{proposition}[Exact domains and equivalences are traces]
\label{prop:ep-exact-trace-surj}
Across any tower prefix, each nonempty trace of an original exact
domain is current exact, and each current exact domain is the trace
of an original exact domain. Equation~\eqref{eq:ep-tp} remains valid
across arbitrary prefix sorts whenever both input traces are nonempty.
For a realised current sort \(p\), restriction
\[
 \rho_p:\operatorname{ExEq}(M)\longrightarrow
                  \operatorname{ExEq}(p),\qquad
 H\longmapsto H|_{p^2}
 \tag{E1}\label{eq:ep-rho}
\]
is a surjective lattice homomorphism.
\end{proposition}
\begin{proof}
For a reachable old word profile \(P\) with a nonempty current
trace, choose endpoints of that trace and realise the word
between them in the old source. Read the actual intermediate
child types from the current homogeneous prefix. Repeated
\eqref{eq:ep-tp}, with these nonempty displayed input traces,
gives the trace of \(P\) as one reachable lifted profile.
Intersections give the assertion for old exact domains. Conversely,
forgetting the child labels of a lifted word, and applying
\eqref{eq:ep-tp}, shows that its endpoint domain is the trace of
the old word profile. Finite nonempty intersections give
surjectivity on exact domains. Iterate through the tower.

Old equivalence traces are reflexive and exact.
Meets commute with set restriction; joins are represented
products of commuting exact equivalences and commute with
restriction by \eqref{eq:ep-tp}, since the input traces contain
the current identity. For surjectivity, express a current exact
equivalence \(F\) on \(p\) as \(D|_{p^2}\) for original exact
\(D\). Reflexivity and original transitivity give
\(\Id\subseteq D\); original symmetry gives \(D=D^\smile\).
The increasing represented powers \(D^n\) stabilise to an
original full exact equivalence \(\langle D\rangle\).
Their traces are \(F^n=F\), by
\eqref{eq:ep-tp} with middle sort \(p\). Hence
\(\rho_p(\langle D\rangle)=F\).
\end{proof}

Define the \emph{canonical largest old lift}
\[
 L_p(F)=
 \bigvee\{H\in\operatorname{ExEq}(M):\rho_p(H)=F\}.
 \tag{E2}\label{eq:ep-old-largest}
\]
Finiteness and join preservation make this well-defined and give
\[
  \rho_p(H)\le F
     \quad\Longleftrightarrow\quad H\le L_p(F).
 \tag{E3}\label{eq:ep-adjunction}
\]
Thus \(L_p\) is a right adjoint and preserves finite meets.
It is important that we choose the \emph{largest} lift, not
an arbitrary exact predecessor of \(F\).

\begin{lemma}[Original bipartite pivot and maximal gate]
\label{lem:ep-old-gate}
If a current restriction \(c|_{p^2}\) is the odd pivot of a
genuine safe step, then the original atom \(c\) is bipartite.
Let \(Y_c<X_c\) be its original even-walk/connectivity interval,
and \(K_p\) the chosen current maximal exact gate. Then
\[
  Y_{\mathrm{current}}=Y_c|_{p^2},\quad
  O_{\mathrm{current}}=O_c|_{p^2},\quad
  X_{\mathrm{current}}=X_c|_{p^2},
 \tag{E4}\label{eq:ep-parity-trace}
\]
and \(L_p(K_p)\) is an \emph{original maximal exact complement}
to \(Y_c<X_c\).
\end{lemma}
\begin{proof}
A bipartite current pivot admits no odd closed \(c\)-walk
inside \(p\). Were the original \(c\)-graph nonbipartite,
some odd \(c^{2r+1}\) would contain the original identity
orbital by transitivity. Repeating (TP) with every
intermediate vertex restricted to \(p\) would put
\(\Id_p\) in \((c|_{p^2})^{2r+1}\), a contradiction.
Thus \(c\) was originally bipartite. Its stabilised even,
odd and connectivity relations are represented old profiles;
repeated (TP) yields \eqref{eq:ep-parity-trace},
including nonemptiness of the odd trace.

Put \(\widehat K=L_p(K_p)\). The adjunction gives
\(Y_c\le\widehat K\). Since \([X_c:Y_c]=2\),
the intersection \(\widehat K\cap X_c\) is either \(Y_c\)
or \(X_c\). The latter would give a nonempty odd trace
inside \(K_p\), contrary to the current complement property.
Thus \(\widehat K\cap X_c=Y_c\).
Any larger old exact complement would have a trace
strictly above \(K_p\), since \(\widehat K\) is the largest
lift, and would contradict current maximality.
\end{proof}

We next isolate the finite support relations on \emph{types}.
Let \(T_j\) be all realised full unary prefix labels, and define
for original invariant \(D\)
\[
   \Phi_j(D)=
   \{(s,t)\in T_j^2:
         \exists x\in s\,\exists y\in t\ (xDy)\},
   \tag{E5}\label{eq:ep-Phi}
\]
using actual equality for \(\Phi_j(\Id)=\Id_{T_j}\).
A loop in \(\Phi_j(a)\) for an old \emph{diversity} atom
means a pair of \emph{distinct} points of the same unary type,
never logical equality.

\begin{lemma}[Finite type supports and exact meets]
\label{lem:ep-phi-support}
Every \(\Phi_j\) preserves union, converse, identity
and represented relational composition. Its restriction to
original exact equivalences also preserves finite meets.
For every original invariant equivalence \(E\),
\(\Phi_j(E)\) is an equivalence on \(T_j\).
An actual \(E\)-class meeting type \(p\) meets
\emph{exactly every} type in the \(\Phi_j(E)\)-class of \(p\).
\end{lemma}
\begin{proof}
Only composition needs an argument. If
\(s\,\Phi_j(D)\,t\) and
\(t\,\Phi_j(F)\,u\), choose actual \(D\)- and
\(F\)-pairs with their middle endpoints both of type \(t\).
By homogeneity of \(M_j\), align those middle points with
an automorphism; the resulting pair has an actual
\(D;F\) witness. The converse inclusion is immediate.
This proves that images of equivalences are equivalences.

For exact \(D_i\) with nonempty original intersection,
nonempty traces of all \(D_i\) on any chosen pair of
current sorts have a common atom by repeated
\cref{lem:ep-exact-meet-prime}. Induction over stages
gives
\[
   \Phi_j\Bigl(\bigcap_iD_i\Bigr)
       =\bigcap_i\Phi_j(D_i)
       \qquad(\bigcap_iD_i\ne\varnothing).
\]
All exact equivalence intersections contain \(\Id\),
so this applies to their meets. Do not extend it to an
empty meet or arbitrary nonexact atomic supports.

If one \(E\)-class \(X\) contains an \(s\)-point and
\((s,t)\in\Phi_j(E)\), use a current automorphism sending
the \(s\)-endpoint of a witnessing \(E\)-pair into \(X\).
Its \(t\)-endpoint stays in \(X\). Conversely types
in \(X\) are \(E\)-related, proving the occupancy assertion.
\end{proof}

\begin{lemma}[One original gate per connected step]
\label{lem:ep-single-old-gate}
For a fixed connected full rigid-groupoid character step,
the canonical largest old lifts \(L_p(K_p)\) of the
active current gates \(K_p\) are \emph{equal} for all
active sorts \(p\). Write their common value \(G\).
Every current full exact bridge in that component is a
nonempty trace of an original full exact \(G\)-quotient
functional correspondence; conversely any such original
correspondence with nonempty trace is a current full bridge.
\end{lemma}
\begin{proof}
For a current exact bridge \(R:p\to q\) with
\(RR^\smile=K_p\), lift its trace to original full exact
\(D\), by \cref{prop:ep-exact-trace-surj}.
The trace product gives
\((D^2)|_{p^2}=K_p\), so the stable original exact
equivalence closure of \(D^2\) lies in
\(G_p=L_p(K_p)\).
Put \(C=G_p;D\). Original one-sort symmetry makes the
old invariant relations commute. Since
\(\Id\subseteq D^2\subseteq G_p\),
we have \(C^2=G_p\); it is an original full exact
functional \(G_p\)-quotient correspondence. Its trace
between \(p,q\) is \(R\), and by TP its square trace
on \(q\) is \(K_q\). Thus \(G_p\) is an old lift of
\(K_q\), whence \(G_p\le G_q\). The reverse direction is
symmetric. Connectedness gives one gate \(G\).

The argument also lifts each bridge to \(C\).
Conversely the nonempty trace of an original \(G\)-bridge
has product with its converse \(G|_{p^2}=K_p\) and is
current exact; fullness of the selected current groupoid
includes it and its target.
\end{proof}

\begin{proposition}[The common exact kernel of the labels]
\label{prop:ep-common-kernel}
Let \(G_1,\ldots,G_m\) be the original gates given by
\cref{lem:ep-single-old-gate}, with repetitions allowed.
Put \(K_{\rm com}=\bigcap_{i=1}^mG_i\) (or \(\Top\)
if \(m=0\)). This is the largest original full exact
equivalence all of whose classes lie in single final
unary type classes. In particular,
\[
  \Phi_m(K_{\rm com})=\Id_T,
  \qquad \bigcap_{i=1}^m\Phi_m(G_i)=\Id_T.
 \tag{E6}\label{eq:ep-common-kernel}
\]
\end{proposition}
\begin{proof}
If \(xK_{\rm com}y\), then their labels agree after
each step: when they are in an active common parent \(p\),
their \(G_i\)-relation lies inside its phase gate \(K_p\),
so their new bits agree; inactive parents are unchanged.

Conversely let original full exact \(E\) preserve
every final label. It preserves every prefix.
At an active parent \(p\) of step \(i\), both
\(E_p=E|_{p^2}\) and \(K_p\) preserve the new bit,
as does their exact represented join. If
\(E_p\not\le K_p\), maximality of the gate would put
the selected odd involution into that strict join,
contradicting preservation of the new bit. Hence
\(E_p\le K_p\) and \eqref{eq:ep-adjunction} gives
\(E\le G_i\). This holds for all steps. Finally use
the nonempty exact-meet preservation in
\cref{lem:ep-phi-support}; the old intersection
contains \(\Id\).
\end{proof}

\begin{theorem}[Exact quotient commutation for safe towers]
\label{thm:ep-quotient-tower}
Let \(E\) be an original full exact equivalence.
Equip \(M/E\) with its actual saturated original quotient
binary orbitals. At prefix \(j\), label an \(E\)-class
by the \(\Phi_j(E)\)-class of the types which occur in it.
These labels form \(T_j/\Phi_j(E)\).
Up to renaming labels, the resulting sequence is itself a
genuine safe maximal-character tower after removing exactly
the steps \(i\) for which \(E\not\le G_i\). A kept step has
canonical old gate \(G_i/E\), after a prescribed
objectwise phase correction. The quotient tower has the
\emph{full} admissible character age, not merely a selected
orbit of markings. If the original tower is terminal,
the quotient tower is terminal.
\end{theorem}
\begin{proof}
We record the details needed for the diversity-rank induction.

\emph{Actual quotient sources.}
For each occupancy label \(C\in T_j/\Phi_j(E)\) choose
a current type \(p_C\in C\).
By \cref{lem:ep-phi-support}, every actual \(E\)-class
of occupancy \(C\) meets \(p_C\), so that quotient sort
is exactly \(p_C/(E|_{p_C^2})\).
The joint exact-axis theorem of Part~I makes these
sorts jointly homogeneous, AP and homomorphic CCO.
Its binary pair orbitals are precisely the restrictions
of original saturated relations:
\[
 (E|_{p_C^2});(a|_{p_C\times p_D});(E|_{p_D^2})
   =(E;a;E)|_{p_C\times p_D},
 \tag{E7}\label{eq:ep-sat-quotient}
\]
by trace-product. This calculation is restricted to realised old atomic
traces \(a|_{p_C\times p_D}\ne\varnothing\); no claim
of nonempty trace is made for an arbitrary old atom.
Conversely, let an old quotient orbital
\(b=E;a;E\) occur between occupancy classes \(C,D\).
Choose representatives of its actual \(E\)-classes
in \(p_C,p_D\). Their old pair has a realised
atomic colour \(a'\), and the quotient orbital
partition implies \(E;a';E=b\).
Applying \eqref{eq:ep-sat-quotient} to this
\emph{realised} \(a'\) shows that every such pair
belongs to the same \(p_C/E_{p_C},p_D/E_{p_D}\)
atomic orbital. Thus old quotient pair colours
together with occupancy unary labels specify
\emph{exactly} the binary orbitals of this
current quotient; there is no hidden finer
binary relation.
In the original one-sort quotient,
an \(E\)-bisaturated full exact old relation descends to
a full exact quotient relation: if
\(D=\bigcap_iP_i=E;D;E\) is exact, use the
exact saturation-intersection lemma to obtain
\(D=\bigcap_i(E;P_i;E)\), where each factor descends
from the old defining word. This supplies the
converse-to-pullback exactness used below.

\emph{A discarded step.}
Suppose \(E\not\le G_i\).
For every active parent \(p\),
\eqref{eq:ep-adjunction} gives \(E_p\not\le K_p\).
Its join \(F_p=E_p\vee K_p\) is a strict
current exact equivalence, hence contains the
\emph{whole} selected phase-one \(K_p\)-quotient
involution relation \(t_p\).
In particular \emph{every actual \(E_p\)-class}
meets both new children, not just some one class:
given any \(x\in p\), choose \(z\) with
\(x\,t_p\,z\). Then \(xF_pz\).
The original same-sort symmetric invariant
equivalences \(E_p,K_p\) commute, so
\(F_p=E_p;K_p\). Consequently there is
\(y\) with \(xE_pyK_pz\). The original
\(K_p\)-gate preserves the new bit, while
\(t_p\) reverses it; hence \(y\) is in
the \(E_p\)-class of \(x\) and has the
opposite bit. This proves the assertion
for every \(x\), without moving
a marked cross-phase pair by an automorphism
of the unmarked parent (which would be invalid).
Thus both children above each active
\(p\) lie in the same new occupancy class.
Occupancy connects all parents of an old
\(\Phi(E)\)-class, and distinct old classes cannot merge.
The quotient unary partition has not changed.

\emph{A retained step.}
Suppose \(E\le G_i\).
Any old occupancy class \(C\) containing one active
parent consists entirely of active parents, because
its restricted old identity \(G_i\)-bridges belong to
the \emph{full} connected current groupoid.
For parents \(p,q\in C\) let
\(I_{pq}=G_i|_{p\times q}\). This is a
\emph{whole} current exact \(K_p,K_q\)-bridge:
\(E\le G_i\) gives a nonempty trace, and
\cref{lem:ep-single-old-gate,lem:ep-phi-support}
give its functional bridge equations and occupancy.
Its character satisfies
\(\chi(I_{pq})+\chi(I_{qr})=\chi(I_{pr})\).
Choose \(p_C\), put \(\gamma_p=\chi(I_{p_Cp})\), and
define on an actual old \(E\)-class \(X\) of occupancy \(C\)
\[
  \zeta(X)=\operatorname{bit}(x)+\gamma_p,
  \quad x\in X\cap p.
 \tag{E8}\label{eq:ep-gamma}
\]
It is well defined: \(E|_{p\times q}\subseteq I_{pq}\)
and \(E|_{p^2}\le K_p\).
The two new occupancy classes are
\[
  C^b=\{p_{b+\gamma_p}:p\in C\},
            \qquad b\in\mathbb F_2 .
 \tag{E9}\label{eq:ep-graph-offsets}
\]
Both occur: the two children above \(p_C\)
are nonempty and cannot meet the same \(E\)-class,
since \(E|_{p_C^2}\le K_{p_C}\) preserves the phase.
Every old \(E\)-class of occupancy \(C\) meets each
parent \(p\in C\), so each refined class projects
bijectively to \(C\). Inactive classes have one lift.

\emph{The full groupoid and the full finite age.}
Identify an active quotient sort \(C\) with
\(p_C/E_{p_C}\), with proposed quotient gate
\(K_{p_C}/E_{p_C}\).
The old \(G_i/E\) is full exact in the
\emph{original one-sort quotient} by the precise
symmetric one-sort exact-descent lemma
\cref{lem:one-original-sort-exact-descent}:
commutation of old invariant orbitals with \(E\)
turns an $E$-saturated full exact old relation
into a quotient word-exact relation.
Its nonempty current trace is the displayed
\emph{full exact} gate by
\cref{prop:ep-exact-trace-surj}, applied to
the \emph{earlier quotient prefix already constructed
by induction on the tower stages}. This invocation
does not presuppose the current, not-yet-retained
character step. (Such downward exact
descent is not asserted for an arbitrary
bisaturated word between independently
quotiented original sorts.)
Now \(E_p\le K_p\) for each active current
parent \(p\). Apply the precise conditional
bridge-descent/lifting lemma
\cref{lem:foundational-rigid-quotient-bridge}
to these actual named parent sorts: current
\(K_p,K_q\)-bridges descend, and \emph{every}
projected full exact bridge pulls back to
an old full exact \(K_p,K_q\)-bridge.
Thus the quotient contains the entire
descended groupoid, not merely a selected
family of transported arrows.
Normalize a bridge \(R:p\to q\) by
\[
  \bar\chi(\bar R)=
           \chi(R)+\gamma_p+\gamma_q .
 \tag{E10}\label{eq:ep-projected-character}
\]
More explicitly, an arrow \(R:p\to q\)
represented instead on \(p_C,p_D\) becomes
\(I_{p_Cp};R;I_{qp_D}\). Its character is
\(\gamma_p+\chi(R)+\gamma_q\), proving the
displayed normalization is independent of
the chosen parent representatives.
This also makes every identity bridge between
two parents of the \emph{same} occupancy class
into the quotient identity of character zero.
Because \(E_p\le K_p\), quotienting by \(E_p\)
does not identify distinct \(K_p\)-classes.
Different rigid arrows already act differently
on those \(K_p\)-classes by freeness, so no two
arrows collapse. This is a character on the
\emph{full} descended groupoid and sends its
selected involution to one.

Old quotient orbital \(E;a;E\) retains its rigid
versus nonrigid status with respect to the \emph{same
transported \(K_p\)-groupoid}: a refinement of a rigid
\(G_i\)-cylinder is rigid with its forced character,
while a nonrigid pair has endpoints in different
rigid-groupoid orbits and no forced phase. Here we
use \(E_p\le K_p\) for every active parent, and
the complete bridge-pullback assertion just proved;
the statement is not a claim that arbitrary
nonrigid correspondences stay nonfunctional after
an unrelated quotient.

We now prove the essential \emph{converse finite
realisation}, without presupposing that arbitrary
fine representatives already have compatible bits.
Let \(P\) be an arbitrary finite \emph{pair-locally
allowed} marked quotient diagram whose underlying
old quotient binary diagram is legal. First realise
its \emph{unmarked} old quotient diagram as actual
\(E\)-classes \(X_v\) in the old homogeneous exact
quotient \(M/E\), using the full exact quotient
source theorem of Part~I. This fixes the entire
old quotient pair diagram at once.

For every vertex \(v\) in an \emph{active} occupancy
class \(C_v\), choose an actual representative
\(x_v\in X_v\cap p_{C_v}\), possible
by \cref{lem:ep-phi-support}. Put
\[
 q_v=[x_v]_{K_{p_{C_v}}},\qquad
 \epsilon_v=\zeta_P(X_v)+\gamma_{p_{C_v}} .
\]
These are actual \emph{current \(K\)-quotient points}
in the rigid-groupoid component, with the desired
unadjusted original bits. For different components we
work componentwise. Inactive occupancy classes have a
single predetermined child label and are not assigned
a \(K_p\) or an \(\epsilon_v\); they impose no phase
seed equation.

Suppose \(q_v,q_w\) lie in the same rigid-groupoid
orbit and let \(R:q_v\mapsto q_w\) be its unique
rigid arrow. Since \(R\) is an \emph{invariant
quotient orbital}, the actual old quotient pair
orbital between \(X_v,X_w\), read in the
representative parent sorts, is rigid and has
character \(\chi(R)\). Pair-local admissibility
of \(P\) under the projected phase rule
\eqref{eq:ep-projected-character} says
\[
 \zeta_P(X_w)+\zeta_P(X_v)
 =\chi(R)+\gamma_{p_{C_v}}+\gamma_{p_{C_w}} .
\]
Consequently
\(\epsilon_w=\epsilon_v+\chi(R)\).
This includes the case \(q_v=q_w\): here
\(R=K_{p_{C_v}}\) is the identity arrow
and \(\epsilon_v=\epsilon_w\).
Thus if several distinct \(E\)-classes lead to
the same \(K\)-quotient point, the proposed
prescription descends to a \emph{function}
on the finite set of distinct \(K\)-quotient
points. It satisfies
\emph{all} pairwise rigid constraints. The
finite phase-extension lemma
\cref{lem:foundational-finite-phase-extension}
extends them to a single \emph{global} groupoid
character marking on all current \(K\)-quotient
points. On inactive sorts retain their unique fixed label.
Pull the marking back to the actual old real
source. Its legality follows directly from the
full character orbit-seed construction, and it
has not changed any old quotient point or old
binary pair orbital. Because \(E_p\le K_p\)
and the identity-bridge corrections \(\gamma_p\)
are coherent, \eqref{eq:ep-gamma} assigns
each \emph{active} \(X_v\) exactly its prescribed
projected bit \(\zeta_P(X_v)\). Inactive \(X_v\)
have their unique unchanged child label.

Thus the quotient of this one globally marked
old source actually contains the given finite
marked diagram \(P\). The argument works for
\emph{every} pair-locally allowed marked quotient
diagram; no independently selected fine edge
colours or SAP of \(M/E\) are required. This is
the full-age assertion needed by the induction.

\emph{Reachability, maximality and terminality.}
At the representative current parent sort \(p\),
let \(c\) be the selected old atomic odd pivot,
\(Y<X\) its even-walk and connectivity equivalences,
and \(K_p\) its maximal exact complement. Put
\(e=E|_{p^2}\le K_p\). The quotient pivot
\(\bar c=(e;c;e)/e\) is ONE quotient atom
by \cref{lem:typed-simultaneous-blocks}, and
is still a diversity atom, since
\(c\subseteq X\setminus Y\) while
\(e\le K_p\) and \(K_p\cap X=Y\).
All current same-sort invariant relations on \(p\)
commute. Thus its represented even and odd powers
are respectively the images of \(e;c^{2r};e\)
and \(e;c^{2r+1};e\). Upon stabilisation the
quotient parity relations are
\[
  \bar Y=(e\vee Y)/e,\qquad
  \bar O=(e;O;e)/e,\qquad
  \bar X=(e\vee X)/e .
  \tag{E11}\label{eq:ep-quotient-parity}
\]
The even and odd pieces remain disjoint:
\(e\vee Y\le K_p\), whereas
\(e;O;e\) is disjoint from \(K_p\)
because \(K_p\cap O=\varnothing\).
Also
\(X\cap(e\vee Y)=Y\) and
\((e\vee X)/(e\vee Y)\) still has
index two by Dedekind transposition
for the commuting equivalences.
Thus the projected pivot is genuinely bipartite
and reachable; its parity interval does not
collapse into equality.

The proposed quotient gate is
\(\bar K=K_p/e\).
The modular law, using \(e\le K_p\), gives
\[
  K_p\cap(e\vee X)=e\vee(K_p\cap X)
                         =e\vee Y.
\]
Therefore \(\bar K\cap\bar X=\bar Y\).
If some larger quotient full exact equivalence
\(\bar K'>\bar K\) were also a complement,
its inverse image \(K'>K_p\) would be
full exact by \cref{lem:typed-simultaneous-blocks},
and \(K'\cap X=Y\) by the displayed
intersection and \(Y\le K_p\).
This contradicts the original maximality of
\(K_p\). Hence the projected complement is maximal.

Finally \(G_i/E\) is an exact relation of the
\emph{original one-sort} quotient, by
\cref{lem:one-original-sort-exact-descent},
and its trace on the already descended
quotient prefix is exactly \(\bar K\).
To verify that it is the \emph{largest original
quotient lift}, take any original quotient
full exact \(H_Q\) with \emph{this same trace}.
Its old inverse image \(H\) is original full exact
and satisfies \(H|_{p^2}=K_p\), since the current
quotient sort is \(p/e\).
The original largest-lift adjunction then gives
\(H\le L_p(K_p)=G_i\), and consequently
\(H_Q\le G_i/E\). Only the previously constructed
quotient prefix is used when taking this trace,
so the argument does not invoke the step
being proved. Thus kept steps are genuinely
maximal safe character steps and the induction
is noncircular.

A proper finite relative-index invariant equivalence
interval in a final quotient sort would pull back
to such an interval in its current representative
parent \(p\), contradicting terminality. This last
argument applies to \emph{all} invariant equivalences,
not just exact ones.
\end{proof}

\begin{remark}[Rank, notation and false shortcuts]
\label{rem:ep-trace-audit}
If \(E>\Id\), the original quotient \(M/E\) has strictly
fewer original diversity orbitals: each quotient
orbital is the image of an old atom, and the diagonal
absorbs some nonidentity old orbital. This is the
strict induction parameter below. An arbitrary
quotient of a Ramsey/EP expansion does \emph{not}
inherit EP automatically; the full-age conclusion
of \cref{thm:ep-quotient-tower} is the needed substitute.
Trace products require \emph{both input traces
nonempty}; \(\Phi\) preserves only nonempty old
exact meets, not empty atomic meets.  Quotient
identities pull back to the whole gate, not raw
logical equality. The present theorem does not
assume that every old \(G_i\)-class meets every
current sort.
\end{remark}
